% TOP_PROBLEM: {"id": 2302044, "title": "Research Problems in Function Theory — Problem 2.44", "category_id": 8, "category": {"id": 8, "name": "analysis", "display_name": "Analysis", "description": "Limits, continuity, calculus, and function theory.", "slug": "analysis", "order_index": 8, "created_at": "2026-07-31T15:26:25.671Z"}, "status": "open"}
% FIRST_POSTED: null
% ENTRY_KIND: statement_only
% Imported from: batch11.tex; UnsolvedMath ID 2302044
% Compile twice: lualatex 2302044.tex
\documentclass[11pt]{article}
\usepackage{amsmath,amssymb,mathtools}
\usepackage{fontspec}
\usepackage{unicode-math}
\setmainfont{Latin Modern Roman}
\setmathfont{Latin Modern Math}
\usepackage{xurl}
\usepackage[unicode,colorlinks=true,linkcolor=blue,urlcolor=blue,pdftitle={UnsolvedMath Problem 2302044}]{hyperref}
\newcommand{\sref}[2]{\hyperlink{source:#1:#2}{[S#2]}}
\newcommand{\cataloguescope}[1]{\paragraph{Mathematical statement.}#1}
\begin{document}
% UnsolvedMath ID 2302044; source catalogue code AMR-022-2044
\setcounter{section}{2302043}
\section{Research Problems in Function Theory — Problem 2.44}\label{2302044}
{\small\sffamily UnsolvedMath ID 2302044 · Analysis}
\subsection{Definitions and mathematical statement}

\paragraph{Shared notation.}$\mathbb N=\{1,2,\ldots\}$, and $\mathbb Z,\mathbb Q,\mathbb R,\mathbb C$ denote the integers, rationals, reals and complex numbers. Any other conventions are specified locally.


For an entire function $f$, write
\[
 M(r,f)=\max_{|z|=r}|f(z)|,\qquad m_0(r,f)=\min_{|z|=r}|f(z)|.
\]
For nonconstant $f$, its order is $\rho(f)=\limsup_{r\to\infty}\log\log M(r,f)/\log r$;
the lower order is defined using $\liminf$ instead. Every constant entire function, including the zero function, has order and lower order zero. In logarithmic modulus ratios, use the extended-real convention $\log0=-\infty$.
For a nonconstant entire $f$ of order $\rho$, define
\[\nu_f(r)=\#\{z:|z|=r,\ |f(z)|=1\}.\]
This counts distinct points. Use $\log0=-\infty$ and $\log(+\infty)=+\infty$ if an exceptional radius has infinitely many such points.

\paragraph{Statement recorded in the supplied source.}
Must
\[\limsup_{r\to\infty}\frac{\log\nu_f(r)}{\log r}=\rho(f)?\] \sref{2302044}{2}
Update 2.44 disproves the proposed identity with $f(z)=e^z$: there are exactly two unit-modulus points on each circle of positive radius, while $\rho(f)=1$. \sref{2302044}{2}
\subsection{Short English statement}
Does the number of unit-modulus points on large circles grow with the same exponent as the entire function's order?
\subsection{Sources}
\begin{description}
\item[\hypertarget{source:2302044:1}{S1}] ULAM.AI and the UnsolvedMath Contributors, \emph{UnsolvedMath: A Curated Collection of Open Mathematics Problems}, record 2302044 (AMR-022-2044). CC BY 4.0. \url{https://huggingface.co/datasets/ulamai/UnsolvedMath}.
\item[\hypertarget{source:2302044:2}{S2}] W. K. Hayman and E. F. Lingham, \emph{Research Problems in Function Theory} (2018), Chapter 2, Problem 2.44 and Update 2.44. \url{https://arxiv.org/abs/1809.07200}.
\end{description}
\label{end:2302044}
\end{document}
